\section*{Limitations}
\label{sec:limitations}

Our study has three main limitations.
First, we evaluate only one language pair (English-Chinese) and causal steering succeeded only in Llama at L13; Gemma L17 failed despite 88.9\% probe accuracy ($n = 171$) and Qwen L14 ($n = 131$) was also negative.
Typologically diverse language pairs and alternative intervention methods are needed to assess generality.
Null steering results in Qwen and Gemma may reflect limitations of single-direction linear steering rather than absence of causal structure; alternative methods such as activation patching \citep{meng2022locating} may reveal non-linear causal pathways that our approach cannot detect.
Second, the effective dataset size (966 to 1,590 samples per model) is constrained by parametric knowledge verification, which retains only questions where the model confidently knows the ground truth and may bias toward well-represented knowledge; additionally, Qwen3-8B serves as both the translation model and an experimental subject, though COMET-based quality control mitigates this concern, and the fact that Llama and Gemma, architecturally independent of the translation model, show the same asymmetry pattern ($\Delta_{\text{asym}} = +7.3$ and $+6.6$~pp, respectively) provides cross-validation that the effect is not an artifact of Qwen-specific translation biases.
Third, all probing results are based on a single train/test split with bootstrap confidence intervals capturing within-split variability; while peak balanced accuracy remains stable across seeds (Llama $\pm$0.12\%, Gemma $\pm$1.32\%, Qwen $\pm$1.90\%), the specific peak layer can shift.
We acknowledge that $k$-fold cross-validation would provide stronger stability guarantees \citep{belinkov2022probing}; given that our core claims rest on cross-model consistency and directional patterns rather than exact accuracy values, the current cross-seed stability analysis provides partial mitigation, and we commit to reporting full 5-fold CV results in the camera-ready version.
This work used AI coding and writing assistants throughout; Appendix~\ref{app:ai-assistants} details the scope of use.
